\scanpage{56}
% archive sheet number: 55
\uline{Remarque} L'hyp. sur $H'$ est certainement
vérifiée si $G_0$ est \deleted{strictement} très rigide
(p. ex. un tore \unsure{affin|forme|affine}, ou un schéma
abélien)

\deleted{\uline{Cor. 5} \inserted{Sous les conditions du cor. 4, i.e.} Supposons $H'$ plat sur $S'$, soit}

\uline{Proposition}

\margin{\rule[-13\baselineskip]{0.35pt}{13\baselineskip}}$S'$ de type fini \deleted{\unsure{réduit et}} \uline{relativement irréductible}
\deleted{\unsure{et}} sur $k$ \inserted{corps} \uline{parfait}, $\deleted{\unsure{H}}G$ schéma en groupes
comm. de type fini sur $k$, $G'=G\times_S S'$,
$H'$ un ss-schéma en groupes fermé
de $G'$ qui est \deleted{\unsure{plat}} \inserted{ouvert (p. ex. plat)} sur $S'$, et \deleted{dont}
la composante connexe ${H'}^{(0)}_{\eta}$ de la fibre
générique est \uline{rigide} (pour ce dernier
point, il suffit que $G_0$ soit \inserted{très} rigide).
Alors (i) il existe un sous-schéma en
groupes $H$ de $G$, \unsure{simple|lisse} sur $k$, tel que
$H\otimes_k \deleted{\unsure{G}}S'$ soit majoré par $H$ et ait $=$
ensemble sous-jacent. \deleted{Sous ces conditions,}
\deleted{\unsure{proposons}} \deleted{il} Considérons \inserted{alors} le morphisme
canonique $u:G\longrightarrow G/H=\deleted{\unsure{G_2 G}}G_1$, et le
morphisme correspondant \deleted{$G'=G\otimes_k S'\longrightarrow \unsure{G/H}$}
$u':G\otimes_k k'\unsure{\longrightarrow|\text{---}}G\otimes_k k'/H\otimes_k k'=\deleted{\unsure{G'_2}}G'_1$. (ii) \deleted{\unsure{Alors}} il
